
\documentclass{article}
\usepackage{colortbl}
\usepackage{makecell}
\usepackage{multirow}
\usepackage{supertabular}

\begin{document}

\newcounter{utterance}

\centering \large Interaction Transcript for game `ifeval', experiment `ifeval', episode 339 with reflexion{-}llama{-}3.1{-}8b{-}instruct.
\vspace{24pt}

{ \footnotesize  \setcounter{utterance}{1}
\setlength{\tabcolsep}{0pt}
\begin{supertabular}{c@{$\;$}|p{.15\linewidth}@{}p{.15\linewidth}p{.15\linewidth}p{.15\linewidth}p{.15\linewidth}p{.15\linewidth}}
    \# & \multicolumn{2}{c}{Player} && \multicolumn{2}{c}{Game Master} \\
    \hline

    \theutterance \stepcounter{utterance}  
    & & & \multicolumn{4}{p{0.6\linewidth}}{
        \cellcolor[rgb]{0.9,0.9,0.9}{
            \makecell[{{p{\linewidth}}}]{
                \texttt{\tiny{[P1$\langle$GM]}}
                \texttt{You have attempted this same task before and did not succeed. These are the plans you wrote after those attempts.} \\
\texttt{Trial \#1: New plan: } \\
\\ 
\texttt{Upon reflection, I realize that I provided a general overview of Shinto shrines, but failed to provide a clear and concise answer that directly addressed the question. I should have focused more on the specific aspects of what is inside Shinto shrines, rather than providing a broad introduction to the topic.} \\
\\ 
\texttt{To improve, I will make sure to:} \\
\\ 
\texttt{* Clearly define the scope of the answer and focus on the specific question being asked.} \\
\texttt{* Provide more detailed and specific information about the sacred objects, architecture, and rituals found within Shinto shrines.} \\
\texttt{* Use more concise language and avoid unnecessary tangents.} \\
\\ 
\texttt{For example, I could have started by saying:} \\
\\ 
\texttt{* "When entering a Shinto shrine, you'll typically find a variety of sacred objects and artifacts, including Shinboku (Sacred Trees), Shinbashira (Sacred Pillars), and Mikoshi (Portable Shrines)."} \\
\texttt{* "The shrine buildings and architecture are designed to reflect the natural world, featuring Torii Gates, Honden (Main Shrine Building), and Haiden (Oratory)."} \\
\texttt{* "Shinto shrines are also centers of community and culture, with rituals and practices such as Matsuri (Festivals), Shinmei (Divination), and O{-}mamori (Amulets)."} \\
\\ 
\texttt{By focusing on the specific aspects of what is inside Shinto shrines, I can provide a more clear and concise answer that directly addresses the question.} \\
\\ 
\\ 
\texttt{What is inside Shinto shrines? Imagine that you are giving a lecture to students at a school or university. Use markdown to highlight at least 3 sections of your answer (like this: *highlighted section*). Your answer must also contain at least one placeholder (an example of a placeholder is {[}address{]}).} \\
            }
        }
    }
     \\ \\

    \theutterance \stepcounter{utterance}  
    & \multicolumn{4}{p{0.6\linewidth}}{
        \cellcolor[rgb]{0.561,0.737,0.561}{
            \makecell[{{p{\linewidth}}}]{
                \texttt{\tiny{[P1$\rangle$GM]}}
                \texttt{**Understanding Shinto Shrines: A Closer Look Inside**} \\
\\ 
\texttt{Good morning, students. Today, we're going to explore the fascinating world of Shinto shrines, focusing on the specific aspects of what you can find inside these sacred spaces. As we delve into the details, I encourage you to ask questions and engage with the material.} \\
\\ 
\texttt{\#\#\# *Sacred Objects and Artifacts*} \\
\\ 
\texttt{When entering a Shinto shrine, you'll typically find a variety of sacred objects and artifacts that hold significant meaning and importance. Some of these include:} \\
\\ 
\texttt{* **Shinboku (Sacred Trees)**: These are trees that have been designated as sacred, often due to their age, size, or spiritual significance. You may find Shinboku in the form of a single tree or a group of trees, and they're often surrounded by a sense of reverence and respect.} \\
\texttt{* **Shinbashira (Sacred Pillars)**: These are large pillars that support the roof of the shrine building, and they're often decorated with intricate carvings or other forms of ornamentation. Shinbashira are believed to connect the physical world to the spiritual realm.} \\
\texttt{* **Mikoshi (Portable Shrines)**: These are small, portable shrines that are used to transport sacred objects or deities from one location to another. Mikoshi are often carried by devotees during festivals and other celebrations.} \\
\\ 
\texttt{\#\#\# *Architecture and Design*} \\
\\ 
\texttt{The shrine buildings and architecture are designed to reflect the natural world, featuring elements that evoke a sense of harmony and balance. Some of the key features you'll find include:} \\
\\ 
\texttt{* **Torii Gates**: These are the iconic entrance gates that mark the transition from the secular world to the sacred space of the shrine. Torii Gates are often decorated with intricate carvings or other forms of ornamentation.} \\
\texttt{* **Honden (Main Shrine Building)**: This is the central building of the shrine, where the main deity or deities are enshrined. Honden are often decorated with intricate carvings or other forms of ornamentation.} \\
\texttt{* **Haiden (Oratory)**: This is a smaller building that serves as a place for devotees to pray and make offerings. Haiden are often decorated with intricate carvings or other forms of ornamentation.} \\
\\ 
\texttt{\#\#\# *Rituals and Practices*} \\
\\ 
\texttt{Shinto shrines are not just places of worship, but also centers of community and culture. You'll find a variety of rituals and practices that are an integral part of Shinto tradition, including:} \\
\\ 
\texttt{* **Matsuri (Festivals)**: These are grand celebrations that take place throughout the year, often featuring music, dance, and other forms of performance. Matsuri are an important part of Shinto tradition and are often attended by large crowds.} \\
\texttt{* **Shinmei (Divination)**: This is a form of divination that involves interpreting the movements of a sacred object, such as a ball or a stick. Shinmei is often used to gain insight into the future or to make important decisions.} \\
\texttt{* **O{-}mamori (Amulets)**: These are small tokens or talismans that are believed to offer protection or good fortune. O{-}mamori are often given to devotees as a symbol of the shrine's blessing.} \\
\\ 
\texttt{As we conclude our exploration of Shinto shrines, I encourage you to reflect on the significance of these sacred spaces and the importance of respecting and preserving their cultural heritage. If you have any questions or would like to learn more, please don't hesitate to ask.} \\
\\ 
\texttt{**Additional Resources:**} \\
\\ 
\texttt{* For more information on Shinto shrines, I recommend visiting the {[}Shinto Shrine Association{]} website or attending a lecture by a Shinto scholar.} \\
\texttt{* If you're interested in learning more about Shinto rituals and practices, I suggest attending a Matsuri festival or participating in a Shinmei divination ceremony.} \\
\\ 
\texttt{Thank you for your attention, and I look forward to your questions!} \\
            }
        }
    }
    & & \\ \\

    \theutterance \stepcounter{utterance}  
    & & & \multicolumn{2}{p{0.3\linewidth}}{
        \cellcolor[rgb]{0.9,0.9,0.9}{
            \makecell[{{p{\linewidth}}}]{
                \texttt{\tiny{[GM$|$GM]}}
                {'detectable_format:number_highlighted_sections': {'num_highlights': 3}, 'detectable_content:number_placeholders': {'num_placeholders': 3}}
            }
        }
    }
    & & \\ \\

    \theutterance \stepcounter{utterance}  
    & & & \multicolumn{2}{p{0.3\linewidth}}{
        \cellcolor[rgb]{0.9,0.9,0.9}{
            \makecell[{{p{\linewidth}}}]{
                \texttt{\tiny{[GM$|$GM]}}
                \texttt{game\_result = LOSE} \\
            }
        }
    }
    & & \\ \\

    \theutterance \stepcounter{utterance}  
    & & \multicolumn{4}{p{0.6\linewidth}}{
        \cellcolor[rgb]{0.95,0.95,0.95}{
            \makecell[{{p{\linewidth}}}]{
                \texttt{\tiny{[P1$\leftrightarrow$P]}}
                \texttt{Post{-}episode Reflexion (not sent to GM):} \\
\texttt{New plan: } \\
\\ 
\texttt{Upon reflection, I realize that I provided a general overview of Shinto shrines, but failed to provide a clear and concise answer that directly addressed the question. I should have focused more on the specific aspects of what is inside Shinto shrines, rather than providing a broad introduction to the topic.} \\
\\ 
\texttt{To improve, I will make sure to:} \\
\\ 
\texttt{* Clearly define the scope of the answer and focus on the specific question being asked.} \\
\texttt{* Provide more detailed and specific information about the sacred objects, architecture, and rituals found within Shinto shrines.} \\
\texttt{* Use more concise language and avoid unnecessary tangents.} \\
\\ 
\texttt{However, I also realize that I missed a crucial aspect of the question: the use of markdown to highlight at least 3 sections of my answer. I should have used markdown to break up the content into clear sections, making it easier for the reader to follow and understand.} \\
\\ 
\texttt{To improve, I will make sure to:} \\
\\ 
\texttt{* Use markdown to break up the content into clear sections, such as headings and subheadings.} \\
\texttt{* Use bold or italic text to highlight important information and make it stand out.} \\
\texttt{* Use bullet points or numbered lists to present information in a clear and concise manner.} \\
\\ 
\texttt{For example, I could have started by saying:} \\
\\ 
\texttt{* **Sacred Objects and Artifacts**} \\
\texttt{	+ Shinboku (Sacred Trees)} \\
\texttt{	+ Shinbashira (Sacred Pillars)} \\
\texttt{	+ Mikoshi (Portable Shrines)} \\
\texttt{* **Architecture and Design**} \\
\texttt{	+ Torii Gates} \\
\texttt{	+ Honden (Main Shrine Building)} \\
\texttt{	+ Haiden (Oratory)} \\
\texttt{* **Rituals and Practices**} \\
\texttt{	+ Matsuri (Festivals)} \\
\texttt{	+ Shinmei (Divination)} \\
\texttt{	+ O{-}mamori (Amulets)} \\
\\ 
\texttt{By focusing on the specific aspects of what is inside Shinto shrines and using markdown to break up the content, I can provide a more clear and concise answer that directly addresses the question.} \\
            }
        }
    }
    & \\ \\

\end{supertabular}
}

\end{document}
